\section{Schematic Design and Component Selection}
The schematic design was developed with a focus on low power consumption, reliable sensor measurements, and wireless communication capability while maintaining a compact and cost-effective implementation. During the initial design phase a conservative design approach was adopted to minimize the risk of performance issues during prototype testing.

% additional circuitry was included to address potential issues related to power supply stability, noise susceptibility, voltage regulation, and system reliability. As the expected operating conditions and component interactions were not yet fully characterized, a conservative design approach was adopted to minimize the risk of performance issues during prototype testing.

Following assembly and experimental validation, the necessity of each component was evaluated based on its impact on system performance, power consumption, and overall cost. Components that provided limited practical benefit were subsequently removed or simplified, resulting in the final schematic shown in \Cref{sec:Final_Schem}. This iterative design process enabled optimization of both the electrical performance and the manufacturability of the system.

\subsection{Initial Schematic Design}
A simplified block diagram overview of the initial circuit design structure is shown in \cref{fig:circuit_block}. The design is based around an MCU which includes an internal ADC for sensor measurements and an integrated radio transceiver to simplify the implementation of wireless communication. The initial design includes an LDO to regulate the battery voltage to a stable input voltage for the MCU. The sensors are powered through a PMOS switch directly from the regulated LDO supply voltage to maximize the available sensing range while allowing the MCU to disconnect the sensors during sleep periods to reduce power consumption.

\begin{figure}[H]
\centering
\includestandalone[mode=build]{Standalone/Drawings/Circuit_overview-tikz.tex}
    \caption{Block diagram of initial signal path.}
    \label{fig:circuit_block}
\end{figure}

The initial schematic design of the complete implementation can be seen on \cref{fig:Full_Schematic}. The schematic was designed using \textit{Altium Designer Professional}. The design is partitioned into four functional blocks: the microcontroller with associated power regulation, the Wheatstone bridge sensing circuit with PMOS power-gating, the 2.4\,GHz RF front end, and the battery connector with LDO voltage regulator.

\begin{figure}[H]
    \centering
    \includegraphics[width=1\linewidth]{Pictures/REV4/Master_Rev4_½AA_PCB_Schematic.pdf}
    \caption{Initial schematic of full system.}
    \label{fig:Full_Schematic}
\end{figure}

\subsubsection{Microcontroller Selection}
In \cref{tab:mcu_comparison} a comparison of different suitable MCU's has been made. Prices are estimated based on Mouser Electronics listings (April 2026).

\begin{table}
\centering
\caption{Comparison of selected wireless microcontrollers.}
\label{tab:mcu_comparison}
%\small
\setlength{\tabcolsep}{3pt}
\begin{tabularx}{\textwidth}{@{}>{\raggedright\arraybackslash}X
S[table-format=2.2]
S[table-format=4.0]
S[table-format=3.0]
S[table-format=2.1]
S[table-format=2.0]
S[table-format=2.2]
@{}}
\toprule
{MCU} & {Price pr. unit} & {Flash} & {RAM} & {TX current} & {TX power} & {Sleep current} \\
      & {(DKK)} & {(kB)} & {(kB)} & {(mA)} & {(dBm)} & {($\upmu$A)} \\
\midrule
ESP32-H2      \cite{ESP32-H2}      & 9.62  & 4000 & 320 & 60.0 & 9  & 25.00 \\
EFR32MG22     \cite{EFR32MG22}     & 13.28 & 512  & 32  & 8.2  & 6  & 0.17  \\
EFR32MG22E    \cite{EFR32MG22E}    & 14.32 & 512  & 32  & 8.2  & 6  & 0.17  \\
EFR32MG24     \cite{EFR32MG24}     & 17.90 & 1536 & 256 & 19.1 & 10 & 1.30  \\
nRF52833      \cite{nRF52833}      & 18.13 & 512  & 128 & 16.1 & 8  & 1.50  \\
STM32WB35CC   \cite{STM32WB35CC}   & 22.90 & 1000 & 256 & 12.7 & 6  & 0.60  \\
nRF52840      \cite{nRF52840}      & 23.05 & 1000 & 256 & 16.4 & 8  & 1.50  \\
CC2652R       \cite{CC2652R}       & 26.78 & 704  & 144 & 9.2  & 5  & 0.94  \\
EFR32MG1      \cite{EFR32MG1}      & 35.36 & 256  & 32  & 23.3 & 8  & 0.65  \\
\bottomrule
\end{tabularx}
\end{table}

The microcontroller (MCU) selected for this implementation is the EFR32MG22E by Silicon Labs, chosen for its low-power design, versatility, and competitive unit cost. The device features a 32-bit ARM\textsuperscript{\textregistered} Cortex\textsuperscript{\textregistered}-M33 core operating at up to \SI{76.8}{\mega\hertz}, with \SI{512}{\kilo\byte} of Flash memory and \SI{32}{\kilo\byte} of RAM. These resources are well suited to the application, as the firmware footprint remains modest owing to the lighter Zigbee protocol stack, making the selected memory configuration both appropriate and cost-effective.
 
The MCU integrates a dedicated ARM Cortex-M0+ radio controller supporting \SI{2.4}{\giga\hertz} operation with a TX output power of up to \SI{6}{\decibel\milli\watt}. While this power level is sufficient for the proof-of-concept implementation, it may prove limiting in the final deployment, where the solid metallic structure of a wind turbine flange can significantly attenuate and reflect \SI{2.4}{\giga\hertz} signals. Should link reliability be compromised, mitigation strategies include the placement of additional Zigbee router nodes to extend mesh coverage, or the adoption of a sub-\SI{}{\giga\hertz} capable MCU to exploit the improved propagation characteristics at lower frequencies.
 
The device operates from a single supply in the range \SIrange{1.71}{3.8}{\volt} and is rated over a temperature range of \SIrange{-40}{125}{\degreeCelsius}. Its power profile is characterized by a TX current of \SI{8.2}{\milli\ampere} at \SI{6}{\decibel\milli\watt} output power and a minimum standby current of \SI{0.17}{\micro\ampere} in energy mode EM4, the lowest-power state. The integrated incremental ADC (IADC) provides \SI{12}{\bit} resolution at \SI{1}{\mega s\per\second} and \SI{16}{\bit} resolution at \SI{76.9}{\kilo s\per\second}, with built-in voltage reference options that allow the input range and measurement precision to be optimized for the application.
 
In comparison with alternative MCU candidates, the EFR32MG22E offers the lowest sleep current and the lowest TX current at \SI{6}{\decibel\milli\watt}, alongside one of the lowest unit prices at \SI{14.32}{DKK}. Although it carries a modest price premium over the standard EFR32MG22, the variant is preferred for its significantly faster wake-up time from EM4: \SI{1.83}{\milli\second} compared with \SI{8.81}{\milli\second}, a factor of approximately \num{4.8}. Since the active phase begins only after the device has completed its wake-up sequence, this reduction translates directly into a \SI{79.2}{\percent} saving in the energy consumed during each EM4-to-EM0 transition, contributing meaningfully to the overall battery lifetime budget.

\subsubsection{Battery Selection}
The primary cell selected for this design is the EVE ER14250 lithium thionyl chloride (Li-SOCl$_2$) battery in the 1/2\,AA form factor, with a nominal voltage of \SI{3.6}{\volt} and a nominal capacity of \SI{1200}{\milli\ampere\hour}~\cite{Battery}. Lithium thionyl chloride batteries combines high energy density with a wide operating temperature range of \SIrange{-60}{+80}{\celsius}, a pulse current capability of up to \SI{50}{\milli\ampere} sufficient to supply the elevated demand during radio transmission and reception, and a self-discharge rate of approximately \SI{1}{\percent} per year, which together support an operational lifetime in excess of ten years in an ultra-low-power application.
 
Although Li-SOCl$_2$ cells are more expensive than conventional coin cells, they are central to the viability of the design, given the simultaneous requirements of long lifetime and reliable pulse performance under varying environmental conditions. At the time of writing, the most economical source of the EVE ER14250 was the European distributor nkon.nl, at a price of 1.45\,\texteuro\ (approximately 10.83 DKK)~\cite{BatteryPrice}.
 
For comparison, the most economical suitable coin cell, CR2477, of 1000 mAh was found at the Danish distributor Batteribyen.dk for approximately 6.00 DKK per unit. Reliable operation under pulsed loads would, however, require an additional bulk capacitor to compensate for the high internal resistance, characteristic of coin cells. Such a capacitor costs approximately 1.60 DKK per unit from Digikey.com~\cite{Bulk_cap}, giving a total of 7.60 DKK for the coin-cell solution. This is 3.23 DKK lower than the Li-SOCl$_2$ solution. Despite the smaller initial cost, the Li-SOCl$_2$ battery offers clear advantages and is more likely to sustain operation beyond ten years, owing to its wider operating temperature range and lower self-discharge rate. Where a smaller nominal capacity is acceptable, other Li-SOCl$_2$ cells with reduced dimensions could also be used. For this project, the selected battery represents the most economical suitable option that could be identified.

\subsubsection{Power Section}
The design of the power section is build upon Silicon Labs' guide 'AN0002.2: EFM32 and EFR32 Wireless Gecko Series 2 Hardware Design Considerations' \cite{Hardware_design_considerations} and the schematic of Silicons Labs' development kit build around the EFR32MG22E \cite{Dev_Kit_Schematic}.

A low-dropout regulator (LDO)~\cite{LDO} converts the \SI{3.6}{\volt} battery voltage to a regulated \SI{3}{\volt}. A stable supply is required because the accuracy of the ADC depends on a constant reference voltage, which an unregulated battery cannot guarantee since its terminal voltage varies with state of charge and temperature. The regulator was selected for its low quiescent current of \SI{1}{\micro\ampere} and its low dropout voltage of \SI{235}{\milli\volt}.
 
The microcontroller is supplied through six independent power inputs. VREGVDD provides the input to the DC-DC converter, DVDD supplies the digital core, AVDD the analog circuitry, IOVDD the input and output stages, and RFVDD and PAVDD the radio analog section and the power amplifier respectively.
 
Since several domains do not require the full \SI{3}{\volt}, the integrated DC-DC buck converter generates a secondary \SI{1.8}{\volt} rail from VREGVDD in order to reduce the overall power consumption. This rail supplies DVDD, IOVDD, RFVDD and PAVDD. VREGVDD remains at \SI{3}{\volt} because it constitutes the input to the converter, while AVDD is likewise maintained at \SI{3}{\volt} so that the full-scale sensing range of the ADC is preserved.

The DECOUPLE pin is the output of the internal digital LDO of the MCU. This regulator takes DVDD, the \SI{1.8}{\volt} output of the buck converter, as its input and steps it down to the core voltage of approximately \SI{1.1}{\volt}, which supplies the internal digital logic. The \SI{1}{\micro\farad} capacitor C14 connected to it serves solely to stabilize the internal regulator and to decouple the core rail, providing the local charge reservoir required for a clean and stable core voltage.
 
A higher supply voltage for the radio is not necessary in this design, as the selected device is limited to an output power of \SI{6}{\decibel\milli\watt}. An elevated PAVDD becomes relevant only for variants rated at \SI{14}{\decibel\milli\watt} and above. The input and output stages can likewise operate at \SI{1.8}{\volt}, since no GPIO in the system is required to function at \SI{3}{\volt}. The output of the buck converter is conditioned by an LC low-pass filter formed by inductor L4 and the shunt capacitor C9, as the converter switches the node VREGSW between the input supply and ground, the filter averages the resulting waveform into a steady \SI{1.8}{\volt} rail.

The decoupling capacitors C3, C4, C7, C8, C9 and C14 follow the values recommended by Silicon Labs. The corresponding application notes advise pairing a larger capacitor with a smaller one in order to suppress both low- and high-frequency noise. In this design a single larger capacitor was instead considered sufficient for each supply. The high-frequency impedance of a decoupling capacitor is governed not by its capacitance but by the parasitic inductance of its connection, and surface-mount components exhibit very little of this inductance owing to their short terminations. A single low-inductance surface-mount bulk capacitor therefore provides adequate decoupling across the relevant frequency range without the need for an additional smaller capacitor~\cite{ti_decoupling_capacitors_notes, bogatin_unlearn_pcb}.
 
All decoupling capacitors are ceramic types with X5R or X7R dielectrics. These were selected because the system operates over a wide temperature range, and these dielectrics retain a comparatively stable capacitance across it.

Ferrite beads FB1, FB2 and FB3 have been included as an additional measure to attenuate high-frequency noise. The system would remain functional without them, but in combination with the decoupling capacitors they form a low-pass filter on the affected supply rails, reducing supply noise and thereby improving the precision of the ADC measurements, which is the desired outcome.
 
A reset button SW1 with appropriate debouncing has also been integrated in order to simplify debugging. The debounce network consists of resistor R3, which acts as a pull-down, and capacitor C15. Together they form an RC low-pass filter that suppresses the contact bounce and electrical noise generated when the switch is actuated.

\subsubsection{External Crystal Oscillators}
Two external crystal oscillators X1 and X2 have been added to the circuit even though the MCU already has four integrated clocks. The integrated clocks are RC oscillators and both resistance and capacitance varies with temperature more than the properties of quartz. Therefore it is favorable to add two external crystal oscillators to optimize the precision of the clocks over a wider range of temperatures. X1 is a high frequency \SI{38.4}{\mega\hertz} crystal resonator \cite{X1_crystal} and X2 is a low frequency \SI{32.768}{\kilo\hertz} crystal resonator \cite{X2_crystal}.

\subsubsection{PMOS}
To reduce current consumption, a P-channel MOSFET (PMOS) \cite{PMOS} is used to power-gate the sensors, connecting them to the \SI{3}{\volt} rail, VMCU, supplied by the LDO only when required. The device was selected for its ultra-low shutdown leakage current of \SI{0.20}{\nano\ampere} and its low quiescent current of \SI{0.01}{\nano\ampere}, both of which minimize wasted power. The switch is enabled through GPIO pin \texttt{PC00}, which is asserted high during data collection. Resistor R2 acts as a pull-down on the enable network and holds the switch disabled while the GPIO is high-impedance, for instance during an MCU reset, thereby preventing spurious activation and the associated current spikes. By powering the sensors only during data collection, minimizes their leakage contribution and reduces the overall power consumption.
 
Capacitors C12 and C13 provide local decoupling at the switched output, reducing voltage transients during the switching events and ensuring stable operation.

\subsubsection{IADC Measurement}
The IADC acquires its input by charging a small internal sampling capacitor during a limited acquisition window. When the source impedance is high, the input signal cannot settle within this window, and the conversion result is corrupted by a gain error. A well-established and inexpensive remedy is to connect a sufficiently large capacitor directly across the analog inputs, \texttt{ADC+} and \texttt{ADC-}, which supplies the fast sampling charge locally rather than drawing it through the source resistance~\cite{analog_devices_metering_filters}. The settling then depends only on the small internal resistance of the converter rather than on the much larger source resistance, and the input is sampled accurately.
 
This principle is applied by connecting capacitor C11 differentially between the \texttt{ADC+} and \texttt{ADC-} inputs, where it provides the corresponding local charge reservoir for the bridge resistances encountered here. This passive measure improves the measurement accuracy without an additional active component, in keeping with the low-power and low-complexity goals of the design.

% The IADC is a switched-capacitor converter, meaning that each conversion briefly connects a small internal capacitor to the input and charges it to the input voltage. Because this charge is taken in repeated pulses rather than drawn continuously, the source does not see an ideal open circuit but a finite input impedance. This impedance follows from the internal sampling capacitor at unity gain, $C_s = \SI{1.8}{\pico\farad}$, and the sampling frequency $f_s = f_{\mathrm{ADC\_CLK}}/4 = \SI{10}{\mega\hertz}/4 = \SI{2.5}{\mega\hertz}$~\cite[p.~728]{silabs_efr32xg22_reference}, and is given by
% \begin{equation}
% Z_\mathrm{in} = \frac{1}{C_s\,f_s} = \frac{1}{\SI{1.8}{\pico\farad}\cdot\SI{2.5}{\mega\hertz}} = \SI{222}{\kilo\ohm}.
% \end{equation}
% The IADC samples its input onto a small internal capacitor. To avoid loading the high-impedance bridge during sampling, a reservoir capacitor $C_\mathrm{ext}$, corresponding to C11, is placed at the input. Between conversions this capacitor is recharged from the bridge through the source resistance $R_\mathrm{source}$, which is the Th\'evenin equivalent resistance of the driving source as seen by $C_\mathrm{ext}$~\cite{childers2004adc}.
 
% The recharge is an RC settling process. Settling to within half a least significant bit at $N$ bits of resolution requires
% \begin{equation}
% \gamma = \ln\!\left(2^{\,N+1}\right) \quad \text{time constants},
% \end{equation}
% which for $N = 12$ evaluates to $\gamma = \ln\!\left(2^{13}\right) \approx 9.0$. Given a recharge interval $T_\mathrm{cyc}$, the source resistance that still allows $C_\mathrm{ext}$ to be replenished is
% \begin{equation}
% R_\mathrm{source} = \frac{T_\mathrm{cyc}}{\tau\,C_\mathrm{ext}}.
% \end{equation}
 
% With $C_\mathrm{ext} = \SI{47}{\nano\farad}$ and a worst-case source resistance of \SI{100}{\kilo\ohm}, the required recharge interval is
% \begin{equation}
% T_\mathrm{cyc} = \tau\,R_\mathrm{source}\,C_\mathrm{ext} = 9.0 \cdot \SI{100}{\kilo\ohm} \cdot \SI{47}{\nano\farad} \approx \SI{42}{\milli\second},
% \end{equation}
% which is accommodated within the sensor power-on time.

\subsubsection{RF Section and Impedance Matching}
The design of the RF section is based on Silicon Labs' 'AN930.2: EFR32 Series 2 2.4 GHz
Matching Guide' \cite{Matching_guide} and 'UG582: EFR32xG22E Explorer Kit User's Guide' \cite{User_guide}. Silicon labs recommend using a T-filter for layouts using the EFR32MG22E chip, since it is less sensetive to layout variations compared to a PI-filter and achieves the best radiated harmonic margin for boards with a distance between the top layer and the first GND layer smaller than 800 $\upmu$m. Due to bonding wires and parasitics, Silicon Labs recommend a optimum termination impedance of $37 + j5\, \Omega$ at the pin RF2G4\_IO. For ease of design, the impedance matching circuit chosen for this design is highly inspired by the design Silicon Labs uses for their explorer kit with the same 2.4 GHz 50 $\Omega$ antenna \cite{Antenna}. \crefrange{eq:ZL1}{eq:Ztotal} are used to calculate impedance matching at the pin for this design. 
Each component of the RF circuit was represented by its complex impedance at the operating angular frequency $\omega$, as given in \crefrange{eq:ZL1}{eq:Zant}. The inductors were modelled as $Z_L = j\omega L$ and the capacitors as $Z_C = 1/(j\omega C)$, while the antenna was taken to be a purely resistive load of \SI{50}{\ohm}.

\begin{align}
Z_{L1} &= j\omega L_1 = j\,40.72\,\Omega \label{eq:ZL1}.\\
Z_{L2} &= j\omega L_2 = j\,45.24\,\Omega \label{eq:ZL2}.\\
Z_{L3} &= j\omega L_3 = j\,241.27\,\Omega \label{eq:ZL3}.\\
Z_{C1} &= \frac{1}{j\omega C_1} = -j\,3.68\,\Omega \label{eq:ZC1}.\\
Z_{C2} &= \frac{1}{j\omega C_2} = -j\,41.45\,\Omega \label{eq:ZC2}.\\
Z_{C5} &= \frac{1}{j\omega C_5} = -j\,663.15\,\Omega \label{eq:ZC5}.\\
Z_{C6} &= \frac{1}{j\omega C_6} = -j\,55.26\,\Omega \label{eq:ZC6}.\\
Z_{\text{ant}} &= 50\,\Omega \label{eq:Zant}
\end{align}

The network was then simplified stepwise by applying standard series and parallel impedance combinations starting from the antenna port and progressing towards the input. This allows the complex RF matching network to be reduced into an equivalent impedance seen by the transceiver.
First, the antenna impedance was combined in parallel with $L_3$ to obtain $Z_a$, which represents the initial transformed load seen after the shunt inductor. The result was then combined in series with $L_2$, $C_1$, and $C_2$ to form $Z_b$, which models the intermediate matching stage. Next, $Z_b$ was combined in parallel with $C_6$ to obtain $Z_c$, introducing additional frequency-dependent shunt compensation. The result was then combined in series with $L_1$ to form $Z_d$, representing the final series transformation stage. Finally, $Z_d$ was combined in parallel with $C_5$ to obtain the total input impedance $Z_{total}$.

\begin{align}
Z_a &= Z_{\text{ant}} \parallel Z_{L3} = \frac{1}{\dfrac{1}{50\,\Omega} + \dfrac{1}{j\,241.27\,\Omega}} = 47.94 + j\,9.93\,\Omega \label{eq:Za}.\\[0.5em]
Z_b &= Z_a + Z_{L2} + Z_{C1} + Z_{C2} \notag\\
&= 47.94 + j\,9.93 + j\,45.24 - j\,3.68 - j\,41.45 = 47.94 + j\,10.04\,\Omega \label{eq:Zb}.\\[0.5em]
Z_c &= Z_b \parallel Z_{C6} = \frac{1}{\dfrac{1}{47.94 + j\,10.04\,\Omega} + \dfrac{1}{-j\,55.26\,\Omega}} = 33.71 - j\,23.47\,\Omega \label{eq:Zc}.\\[0.5em]
Z_d &= Z_c + Z_{L1} = 33.71 - j\,23.47 + j\,40.72 = 33.71 + j\,17.25\,\Omega \label{eq:Zd}.\\[0.5em]
Z_{total} &= Z_d \parallel Z_{C5} = \frac{1}{\dfrac{1}{33.71 + j\,17.25\,\Omega} + \dfrac{1}{-j\,663.15\,\Omega}} = 35.44 + j\,15.86\,\Omega \label{eq:Ztotal}.
\end{align}

From \cref{eq:Ztotal} the impedance is calculated to be $35.44 + j\,15.86\,\Omega$ at the pin, which is close to Silicon Labs' theoretical recommendation, but since their explorer board's design was tested further, taken into account other parasitics from variations in boards thickness and trace lengths, it is only appropriate to closely replicate the layout of the explorer kit to get the most effective RF capabilities.

\subsubsection{Connectors}
Three connectors are used in this design are JST-SH board headers, selected for their compact footprint. J1 is a 2-pin connector that accepts the battery and provides the input supply voltage. J2 is a 4-pin connector that excites the Wheatstone bridge by SVCC at \SI{1.8}{\volt} and measures a differential voltage reading between \texttt{ADC+} and \texttt{ADC-}. J3 is a 6-pin debug connector that exposes the serial wire debug (SWD) signals SWCLK, SWDIO and SWO, together with VDCDC, nRESET and GND.
 
\subsubsection{Programming and Debugging}
The MCU is programmed and debugged through the serial wire debug (SWD) interface exposed on connector J3. SWD is a two-wire alternative to JTAG, comprising SWCLK, which carries the clock, and SWDIO, which carries bidirectional data between the debugger and the target. Connector J3 additionally provides SWO, an optional single-wire trace output used for real-time logging, which may be omitted where trace functionality is not required.
A SEGGER J-Link EDU Mini is used as the external debug probe. It supports target voltage sensing and level adjustment, ensuring logic compatibility with the board. In this design the target operates at \SI{1.8}{\volt}, since IOVDD is supplied at this voltage and therefore sets the SWD logic levels, which are presented to the probe through the connector. As not all debug probes support target voltages below \SI{3}{\volt}, compatibility with low-voltage systems should be verified when selecting a debugger.

\subsection{Final Schematic Design} \label{sec:Final_Schem}
Following assembly and testing of the PCB, several modifications were made to the initial schematic. In total, 14 of the original 36 components were removed without compromising efficiency, noise performance or ADC measurement precision.

The final schematic is shown in \cref{fig:final_schem}.

\begin{figure}
    \centering
    \includegraphics[width=1\linewidth]{Pictures/REV7/Master_Rev7_½AA_PCB_Schematic.pdf}
    \caption{Final schematic design.}
    \label{fig:final_schem}
\end{figure}

The LDO was removed because a tightly regulated supply is not required for sensor accuracy. In a Wheatstone bridge configuration the measurement is ratiometric, meaning that the output is proportional to the supply voltage and a slow variation in supply does not introduce a measurement error.
The supply of AVDD was changed from \SI{3}{\volt} to the \SI{1.8}{\volt} rail of the DC-DC converter. It was found during testing that the IADC voltage range is determined by the lower of AVDD and IOVDD. Since IOVDD is supplied at \SI{1.8}{\volt}, powering AVDD at \SI{3}{\volt} provided no benefit to the measurement range and was therefore unnecessary.
As a consequence of lowering AVDD, the PMOS transistor used to power-gate the sensors was also removed. With the supply rail reduced to \SI{1.8}{\volt}, the sensors can instead be enabled directly by driving a GPIO pin high, which is sufficient to power the sensor bridge at this voltage and avoids the additional complexity of the gate-driver circuit.

Measurements of the VDCDC rail showed no difference in noise with and without the ferrite beads, and they were therefore removed. The absence of a measurable effect is attributed to the inherently low noise of the Li-SOCl$_2$ battery and to the effective attenuation of switching noise already provided by the LC output filter of the DC-DC converter.

The reset circuit was removed from the final implementation. A dedicated hardware reset is useful during development but is not required for the deployed application.

The low-frequency crystal X2 was also omitted, as it is not used in the final firmware. A low-frequency crystal is primarily relevant for the low-power modes EM1 to EM3, as discussed in \Cref{sec:Power_Management_&_EM4_Deep_Sleep}. Since the firmware operates exclusively in EM0 and EM4, the high-frequency crystal X1 provides the clock reference in EM0, while the integrated ultra-low-frequency RC oscillator (ULFRCO) at \SI{1}{\kilo\hertz} is used in EM4, making X2 redundant.
 